% GENERATED FILE. Edit canonical lessons, then run npm run generate:books.
% canonical-source-hash: fnv1a64:c8d75e96e19fed37
% canonical-lessons: ML-C226-carcca-ceyyuka, ML-C226-visadikarikkuka, ML-C226-tayyarakkuka, ML-C226-purttiyakkuka, ML-C226-parisodhikkuka

\chapter{To discuss, To explain, To prepare, To finish, To check}
\label{ch:ml-to-discuss-to-explain-to-prepare-to-finish-to-check}
\hlchaptermodality{226}

\begin{chapteropening}
\textbf{By the end of this chapter:} \emph{I can say to discuss, to explain, to prepare, to finish and to check.}

Complete the last lesson of chapter 226: I can say to discuss, to explain, to prepare, to finish and to check.
\end{chapteropening}

\section[\texorpdfstring{\ml{ചർച്ച} \ml{ചെയ്യുക} (carcca ceyyuka)}{ചർച്ച ചെയ്യുക (carcca ceyyuka)}]{\ml{ചർച്ച} \ml{ചെയ്യുക} (carcca ceyyuka) — to discuss}

\label{lesson:ML-C226-carcca-ceyyuka}



\begin{quote}
Before the new one: say the Malayalam for to return, then the Malayalam for to take.
\end{quote}

\subsection*{You'll want to know: \ml{ചർച്ച} \ml{ചെയ്യുക}}
\textbf{\ml{ചർച്ച} \ml{ചെയ്യുക}} — \emph{carcca ceyyuka} — \textquotedblleft{}to discuss\textquotedblright{}.

\begin{cousinweb}
\textbf{\ml{ചർച്ച} \ml{ചെയ്യുക}} (\emph{carcca ceyyuka}) joins Sanskrit \emph{carcā}, \textquotedblleft{}repetition, going over\textquotedblright{} (CDIAL 4695), to native \emph{ceyyuka}, \textquotedblleft{}to do\textquotedblright{} (DEDR 1957).
\end{cousinweb}

\begin{grammarlens}[title={What you've built}]
One more word for this chapter.
\end{grammarlens}

\subsection*{Your turn}
\begin{itemize}
\raggedright
  \item \emph{Say it:} \emph{carcca ceyyuka}
  \item \emph{Say it:} \emph{carcca ceyyuka}, once more
  \item \emph{Say it:} say \emph{carcca ceyyuka}
  \item \emph{Recall:} say the Malayalam for to return, then the Malayalam for to take, then say \emph{carcca ceyyuka} again
\end{itemize}

\begin{tcolorbox}[breakable,colback=teal!4,colframe=teal!35!black,title={Before you move on}]
What does \ml{ചർച്ച} \ml{ചെയ്യുക} mean? (\textquotedblleft{}To discuss\textquotedblright{}.) Say it once more.
\end{tcolorbox}
\section[\texorpdfstring{\ml{വിശദീകരിക്കുക} (viśadīkarikkuka)}{വിശദീകരിക്കുക (viśadīkarikkuka)}]{\ml{വിശദീകരിക്കുക} (viśadīkarikkuka) — to explain}

\label{lesson:ML-C226-visadikarikkuka}



\begin{quote}
Before the new one: say the Malayalam for to take, then the Malayalam for to discuss.
\end{quote}

\subsection*{You'll want to know: \ml{വിശദീകരിക്കുക}}
\textbf{\ml{വിശദീകരിക്കുക}} — \emph{viśadīkarikkuka} — \textquotedblleft{}to explain\textquotedblright{}.

\begin{cousinweb}
\textbf{\ml{വിശദീകരിക്കുക}} (\emph{viśadīkarikkuka}) is built on Sanskrit \emph{viśada}, \textquotedblleft{}clear, bright\textquotedblright{}, with \emph{-karikkuka}, \textquotedblleft{}to make\textquotedblright{}, on the Sanskrit root \emph{kṛ}: \textquotedblleft{}to make clear\textquotedblright{}.
\end{cousinweb}

\begin{grammarlens}[title={What you've built}]
One more word for this chapter.
\end{grammarlens}

\subsection*{Your turn}
\begin{itemize}
\raggedright
  \item \emph{Say it:} \emph{viśadīkarikkuka}
  \item \emph{Say it:} \emph{viśadīkarikkuka}, once more
  \item \emph{Say it:} say \emph{viśadīkarikkuka}
  \item \emph{Recall:} say the Malayalam for to take, then the Malayalam for to discuss, then say \emph{viśadīkarikkuka} again
\end{itemize}

\begin{tcolorbox}[breakable,colback=teal!4,colframe=teal!35!black,title={Before you move on}]
What does \ml{വിശദീകരിക്കുക} mean? (\textquotedblleft{}To explain\textquotedblright{}.) Say it once more.
\end{tcolorbox}
\section[\texorpdfstring{\ml{തയ്യാറാക്കുക} (tayyāṟākkuka)}{തയ്യാറാക്കുക (tayyāṟākkuka)}]{\ml{തയ്യാറാക്കുക} (tayyāṟākkuka) — to prepare}

\label{lesson:ML-C226-tayyarakkuka}



\begin{quote}
Before the new one: say the Malayalam for to discuss, then the Malayalam for to explain.
\end{quote}

\subsection*{You'll want to know: \ml{തയ്യാറാക്കുക}}
\textbf{\ml{തയ്യാറാക്കുക}} — \emph{tayyāṟākkuka} — \textquotedblleft{}to prepare\textquotedblright{}.

\begin{cousinweb}
\textbf{\ml{തയ്യാറാക്കുക}} (\emph{tayyāṟākkuka}) joins \emph{tayyāṟŭ}, \textquotedblleft{}ready\textquotedblright{}, from Hindustani \emph{taiyār}, an Arabic word, to native \emph{ākkuka}, \textquotedblleft{}to make\textquotedblright{} (DEDR 333). DEDR 333 compares Sanskrit \emph{ām}.
\end{cousinweb}

\begin{grammarlens}[title={What you've built}]
One more word for this chapter.
\end{grammarlens}

\subsection*{Your turn}
\begin{itemize}
\raggedright
  \item \emph{Say it:} \emph{tayyāṟākkuka}
  \item \emph{Say it:} \emph{tayyāṟākkuka}, once more
  \item \emph{Say it:} say \emph{tayyāṟākkuka}
  \item \emph{Recall:} say the Malayalam for to discuss, then the Malayalam for to explain, then say \emph{tayyāṟākkuka} again
\end{itemize}

\begin{tcolorbox}[breakable,colback=teal!4,colframe=teal!35!black,title={Before you move on}]
What does \ml{തയ്യാറാക്കുക} mean? (\textquotedblleft{}To prepare\textquotedblright{}.) Say it once more.
\end{tcolorbox}
\section[\texorpdfstring{\ml{പൂർത്തിയാക്കുക} (pūrttiyākkuka)}{പൂർത്തിയാക്കുക (pūrttiyākkuka)}]{\ml{പൂർത്തിയാക്കുക} (pūrttiyākkuka) — to finish, to complete}

\label{lesson:ML-C226-purttiyakkuka}



\begin{quote}
Before the new one: say the Malayalam for to explain, then the Malayalam for to prepare.
\end{quote}

\subsection*{You'll want to know: \ml{പൂർത്തിയാക്കുക}}
\textbf{\ml{പൂർത്തിയാക്കുക}} — \emph{pūrttiyākkuka} — \textquotedblleft{}to finish, to complete\textquotedblright{}.

\begin{cousinweb}
\textbf{\ml{പൂർത്തിയാക്കുക}} (\emph{pūrttiyākkuka}) joins Sanskrit \emph{pūrti}, \textquotedblleft{}completion, fulfilment\textquotedblright{}, to native \emph{ākkuka}, \textquotedblleft{}to make\textquotedblright{} (DEDR 333). DEDR 333 compares Sanskrit \emph{ām}.
\end{cousinweb}

\begin{grammarlens}[title={What you've built}]
One more word for this chapter.
\end{grammarlens}

\subsection*{Your turn}
\begin{itemize}
\raggedright
  \item \emph{Say it:} \emph{pūrttiyākkuka}
  \item \emph{Say it:} \emph{pūrttiyākkuka}, once more
  \item \emph{Say it:} say \emph{pūrttiyākkuka}
  \item \emph{Recall:} say the Malayalam for to explain, then the Malayalam for to prepare, then say \emph{pūrttiyākkuka} again
\end{itemize}

\begin{tcolorbox}[breakable,colback=teal!4,colframe=teal!35!black,title={Before you move on}]
What does \ml{പൂർത്തിയാക്കുക} mean? (\textquotedblleft{}To finish, to complete\textquotedblright{}.) Say it once more.
\end{tcolorbox}
\section[\texorpdfstring{\ml{പരിശോധിക്കുക} (pariśōdhikkuka)}{പരിശോധിക്കുക (pariśōdhikkuka)}]{\ml{പരിശോധിക്കുക} (pariśōdhikkuka) — to check, to inspect}

\label{lesson:ML-C226-parisodhikkuka}



\begin{quote}
Before the new one: say the Malayalam for to prepare, then the Malayalam for to finish.
\end{quote}

\subsection*{You'll want to know: \ml{പരിശോധിക്കുക}}
\textbf{\ml{പരിശോധിക്കുക}} — \emph{pariśōdhikkuka} — \textquotedblleft{}to check, to inspect\textquotedblright{}.

\begin{cousinweb}
\textbf{\ml{പരിശോധിക്കുക}} (\emph{pariśōdhikkuka}) is built on Sanskrit \emph{pariśōdhana}, \textquotedblleft{}thorough cleansing, examining\textquotedblright{}, with the ending that makes a Malayalam verb.
\end{cousinweb}

\begin{grammarlens}[title={What you've built}]
One more word for this chapter.
\end{grammarlens}

\subsection*{Your turn}
\begin{itemize}
\raggedright
  \item \emph{Say it:} \emph{pariśōdhikkuka}
  \item \emph{Say it:} \emph{pariśōdhikkuka}, once more
  \item \emph{Say it:} say \emph{pariśōdhikkuka}
  \item \emph{Recall:} say the Malayalam for to prepare, then the Malayalam for to finish, then say \emph{pariśōdhikkuka} again
\end{itemize}

\begin{tcolorbox}[breakable,colback=teal!4,colframe=teal!35!black,title={Before you move on}]
What does \ml{പരിശോധിക്കുക} mean? (\textquotedblleft{}To check, to inspect\textquotedblright{}.) Say it once more.
\end{tcolorbox}
